\documentclass[11pt]{article}
\usepackage{coursenotes}
\begin{document}
\handoutheader{H1}{Instrumental Variables}{Encouragement, noncompliance, and outside shifters of treatment}{Meeting 7 (Observational data)}

\begin{decision}
A firm can randomise an \emph{offer} but not whether people take it up, or it suspects an unmeasured confounder that no
adjustment can fix. An instrument, something that shifts treatment but affects the outcome only through treatment, can
still identify an effect: the effect on the units whose treatment the instrument moved. This handout says when, for whom,
and how to tell a strong instrument from a weak one.
\end{decision}

\begin{goals}
  \item state the four IV assumptions and say which can be checked;
  \item show what randomisation of the instrument alone identifies (the ITT and the first stage);
  \item prove the LATE theorem and recover complier means;
  \item quantify what failures of exclusion and monotonicity do to the ratio;
  \item interpret two-stage least squares with covariates, and guard against weak instruments.
\end{goals}

%=====================================================================================
\sectionone

\subsection{Setup}

Let $Z \in \{0,1\}$ be an instrument, $D \in \{0,1\}$ a treatment and $Y$ an outcome. Each unit has \emph{potential treatments}
$D(0), D(1)$ and potential outcomes $Y(d, z)$; we observe $D = D(Z)$ and $Y = Y(D, Z)$.

\begin{assum}{Independence}{ind}
$\{Y(d,z), D(z) : d, z \in \{0,1\}\} \indep Z$.
\end{assum}
\begin{assum}{Exclusion}{excl}
$Y(d, 0) = Y(d, 1) \equiv Y(d)$: the instrument affects the outcome only through treatment.
\end{assum}
\begin{assum}{Relevance}{rel}
$\E[D(1) - D(0)] \ne 0$.
\end{assum}
\begin{assum}{Monotonicity}{mono}
$D(1) \ge D(0)$ for every unit.
\end{assum}

A randomised instrument delivers independence by design. Relevance is checkable. Monotonicity is a claim about behaviour,
argued from the institution. Exclusion is a claim that no other channel exists, and in general it can be tested neither from
the data nor from the design.

\subsection{What independence alone buys}

\begin{prop}{Two effects of the instrument}{itt}
Under independence alone, $\text{ITT} \equiv \E[Y \mid Z=1] - \E[Y \mid Z=0] = \E[Y(D(1),1) - Y(D(0),0)]$ and
$\pi \equiv \E[D \mid Z=1] - \E[D \mid Z=0] = \E[D(1) - D(0)]$.
\end{prop}
\begin{proof}
Given $Z = z$, $Y = Y(D(z), z)$ and $D = D(z)$; independence lets the conditioning be dropped.
\end{proof}

The ITT is the effect of assigning the instrument through every channel. If the business controls the offer but not its
take-up, and would roll the offer out with everything that came with it in the test, the ITT \emph{is} the policy effect and
exclusion is not needed. Exclusion is needed only to attribute the effect to the treatment itself.

\subsection{Compliance types and the LATE}

Under monotonicity each unit is an \emph{always-taker} ($D(0) = D(1) = 1$), a \emph{never-taker} ($D(0) = D(1) = 0$) or a
\emph{complier} ($D(0) = 0$, $D(1) = 1$); \emph{defiers} ($D(0) = 1$, $D(1) = 0$) are ruled out.

\begin{prop}{Type shares}{shares}
Under independence and monotonicity, $\pi_A = \Pr(D=1 \mid Z=0)$, $\pi_N = \Pr(D=0 \mid Z=1)$ and $\pi_C = \pi$.
\end{prop}
\begin{proof}
With no defiers, $D = 1$ under $Z = 0$ only for always-takers and $D = 0$ under $Z = 1$ only for never-takers; independence makes
each arm's shares the population's. The rest are compliers, and $\pi = (\pi_A + \pi_C) - \pi_A$.
\end{proof}

No individual's type is observed; only shares and averages over types are identified.

\begin{prop}{The LATE theorem (Imbens and Angrist)}{late}
Under Assumptions~\ref{asm:ind}--\ref{asm:mono},
\[
  \frac{\text{ITT}}{\pi} = \frac{\E[Y \mid Z=1] - \E[Y \mid Z=0]}{\E[D \mid Z=1] - \E[D \mid Z=0]} = \E[Y(1) - Y(0) \mid \text{complier}] \equiv \text{LATE}.
\]
\end{prop}
\begin{proof}
By exclusion the ITT is $\E[Y(D(1)) - Y(D(0))]$. For always- and never-takers $D(1) = D(0)$ and the difference is zero; for compliers
it is $Y(1) - Y(0)$; there are no defiers. So $\text{ITT} = \pi_C\,\text{LATE}$; divide by $\pi = \pi_C > 0$.
\end{proof}

The ratio (the \emph{Wald estimator}) is the effect on units whose treatment the instrument changed. The data say nothing
about always-takers or never-takers.

\begin{prop}{Complier outcome means}{complier}
Under Assumptions~\ref{asm:ind}--\ref{asm:mono},
$\E[Y(1) \mid C] = \{\E[YD \mid Z=1] - \E[YD \mid Z=0]\}/\pi$ and $\E[Y(0) \mid C] = \{\E[Y(1-D) \mid Z=0] - \E[Y(1-D) \mid Z=1]\}/\pi$.
\end{prop}
\begin{proof}
$\E[YD \mid Z=1] = \pi_A\E[Y(1) \mid A] + \pi_C\E[Y(1) \mid C]$ and $\E[YD \mid Z=0] = \pi_A\E[Y(1) \mid A]$; subtract and divide. The second uses
never-takers in the same way.
\end{proof}

Replacing $Y$ with a covariate describes the compliers, the first step in arguing whether their effect transports.

\subsection{When exclusion or monotonicity fails}

\begin{prop}{Exclusion violation}{violation}
Keep independence, relevance and monotonicity, but let $Y(d,1) = Y(d,0) + \delta$ for every unit. Then
$\text{ITT}/\pi = \text{LATE} + \delta/\pi$.
\end{prop}
\begin{proof}
$Y(D(1),1) - Y(D(0),0) = \{Y(D(1),0) - Y(D(0),0)\} + \delta$; the first part averages to $\pi\,\text{LATE}$.
\end{proof}

The direct effect is divided by the first stage: harmless with a strong instrument, fatal with a weak one. Report the LATE as a
function of $\delta$ and the $\delta$ at which the decision flips.

\begin{prop}{Defiers}{defiers}
Keep independence, exclusion and relevance but allow defiers. Then
$\text{ITT}/\pi = (\pi_C\,\text{LATE}_C - \pi_D\,\text{LATE}_D)/(\pi_C - \pi_D)$.
\end{prop}
\begin{proof}
The LATE proof now leaves a defier term $-\pi_D\,\text{LATE}_D$ in the ITT, and the first stage is $\pi_C - \pi_D$.
\end{proof}

With defiers one weight is negative, and the ratio need not lie between the two groups' effects; it can have the wrong sign.

\subsection{Covariates and two-stage least squares}

If the instrument is random only given covariates $X$ (assignment probabilities differ by stratum), two-stage least squares with
stratum indicators computes $\hat\theta = \sum_i\tilde Z_iY_i/\sum_i\tilde Z_iD_i$ with $\tilde Z_i = Z_i - \hat\E[Z \mid X_i]$.

\begin{prop}{What 2SLS with strata estimates}{tsls}
Under the conditional versions of Assumptions~\ref{asm:ind}--\ref{asm:mono}, with discrete $X$,
\[
  \hat\theta \to \frac{\E[\Var(Z \mid X)\,\pi(X)\,\text{LATE}(X)]}{\E[\Var(Z \mid X)\,\pi(X)]}.
\]
\end{prop}
\begin{proof}
For binary $Z$, $\E[\tilde ZY] = \E[\Var(Z \mid X)\,\text{ITT}(X)]$ and $\E[\tilde ZD] = \E[\Var(Z \mid X)\,\pi(X)]$; within each stratum
$\text{ITT}(X) = \pi(X)\,\text{LATE}(X)$.
\end{proof}

Controls that determine the instrument's probability are required; controls that only predict $Y$ reduce variance; anything
measured after assignment is never a control. With heterogeneous effects, 2SLS weights strata towards balanced instruments and
strong first stages. Running the two stages by hand reproduces the estimate but not its standard error; use an IV routine. With a
continuous treatment (a log price) the same formula gives an elasticity, a weighted average of slopes over units whose treatment
the instrument moved.

\subsection{Inference and weak instruments}

\begin{prop}{Delta-method standard error}{delta}
With $\hat\theta = \widehat{\text{ITT}}/\hat\pi$, approximately
$\Var(\hat\theta) \approx \{\Var(\widehat{\text{ITT}}) - 2\theta\Cov(\widehat{\text{ITT}}, \hat\pi) + \theta^2\Var(\hat\pi)\}/\pi^2$.
\end{prop}
\begin{proof}
Linearise $a/b$ at $(\text{ITT}, \pi)$: $a/b - \theta \approx \{(a - \text{ITT}) - \theta(b - \pi)\}/\pi$.
\end{proof}

With a precise first stage, $\SE(\hat\theta) \approx \SE(\widehat{\text{ITT}})/\pi$. The approximation fails when $\pi$ is small relative to its
SE, measured by the first-stage $F = (\hat\pi/\SE(\hat\pi))^2$. A weak instrument's estimate is heavy-tailed, centred between the truth
and the naive regression, and its nominal 95\% interval undercovers. $F > 10$ (Staiger and Stock) is a floor, not a certificate: Lee,
McCrary, Moreira and Porter (2022) show a conventional 5\% $t$-test has correct size only for $F$ above about $104.7$.

\begin{prop}{The Anderson--Rubin test}{ar}
If $Y = \theta D + u$ with $Z \indep u$, the regression of $Y - \theta_0D$ on $Z$ has slope zero exactly when $\theta_0 = \theta$ (given $\pi \ne 0$),
and the usual test of that slope has correct size whatever the first stage's strength.
\end{prop}
\begin{proof}
$Y - \theta_0D = u + (\theta - \theta_0)D$, whose covariance with $Z$ is $(\theta - \theta_0)\Cov(Z, D)$; under the null the dependent variable is
$u$, independent of $Z$.
\end{proof}

The set of $\theta_0$ not rejected is a confidence set valid under weak instruments; it may be very wide or unbounded, and that
width is an accurate summary of what the data can say. Decide the reporting rule before computing the estimate: report the first stage and $F$; below 10, report the ITT and first stage
rather than the ratio; between 10 and about 100, report the Anderson--Rubin set.

\subsection{Locality}

The LATE averages over compliers, defined by the instrument; a different instrument has different compliers and, with heterogeneous
effects, a different LATE. Ask who the compliers are (Result~\ref{prop:complier}), how large a change in treatment the instrument
induced, and whether the decision is about the compliers at all.

\begin{pitfall}[IV errors]
Comparing those who took the treatment with those who did not; dividing a small ITT by a weak first stage and reporting the ratio;
controlling for anything measured after assignment; reading a LATE as the ATE; ignoring a direct effect of the encouragement.
\end{pitfall}

%=====================================================================================
\sectiontwo

\paragraph{Setting.} QuickBite emails a random half of 20{,}000 customers an invitation to join QuickBite Plus ($Z = 1$). Joining ($D$)
is up to the customer. The outcome $Y$ is orders over the next month. Each order earns ¥10 of margin; each Plus member costs ¥15 a month
in waived delivery fees.

\paragraph{Step 1: ITT and first stage.} Plus membership is $0.30$ among invited and $0.05$ among others (customers can join on their own),
so $\pi = 0.25$, $\SE(\hat\pi) = 0.0051$ and $F \approx 2{,}400$: a very strong instrument. Mean orders are $4.30$ and $4.00$, so
$\text{ITT} = 0.30$ (SE about $0.05$, with an outcome SD of 3.5).

\paragraph{Step 2: types and the LATE.} $\pi_A = 0.05$, $\pi_N = 0.70$, $\pi_C = 0.25$. $\text{LATE} = 0.30/0.25 = 1.2$ extra orders a month per
complier, SE about $0.05/0.25 = 0.20$. Complier means (Result~\ref{prop:complier}): with $\E[YD \mid Z=1] = 1.75$ and $\E[YD \mid Z=0] = 0.35$,
$\E[Y(1) \mid C] = 1.40/0.25 = 5.6$; with $\E[Y(1-D) \mid Z=0] = 3.65$ and $\E[Y(1-D) \mid Z=1] = 2.55$, $\E[Y(0) \mid C] = 4.4$. Compliers would order
$4.4$ times a month without Plus and $5.6$ with it.

\paragraph{Step 3: the decisions.} Two different questions.
\begin{itemize}
  \item \emph{Should QuickBite send the invitation?} The ITT is the policy effect: per invited customer, $10 \times 0.30 = $ ¥3.00 of margin against
  $0.25 \times 15 = $ ¥3.75 of waived fees for the extra members. Net $-$¥0.75: do not send it, whether or not exclusion holds.
  \item \emph{Is Plus worth it for the customers the email recruits?} $10 \times 1.2 = $ ¥12 a month against ¥15: no, for compliers. Always-takers
  are a different population; the email says nothing about them.
\end{itemize}

\paragraph{Step 4: exclusion.} The email itself may remind customers to order. If it adds $\delta = 0.05$ orders directly, the ratio
overstates the LATE by $0.05/0.25 = 0.2$ (Result~\ref{prop:violation}): the true complier effect would be $1.0$, strengthening the conclusion.
A reminder email without the invitation, sent to a third randomised group, would measure $\delta$.

\begin{exercise}[computation]
Suppose instead the invitation lifted membership only from $0.05$ to $0.07$ ($\pi = 0.02$, $\SE(\hat\pi) \approx 0.0034$) with the same ITT of $0.30$
(SE $0.05$). Compute the first-stage $F$, the Wald ratio and its approximate SE, and say what you would report.
\end{exercise}
\answer{$F = (0.02/0.0034)^2 \approx 35$. Ratio $0.30/0.02 = 15$ orders per complier, SE roughly $0.05/0.02 = 2.5$ before the first-stage term,
and more after it. An effect of 15 orders a month per complier is implausible: any direct effect of the email of $\delta$ is multiplied by
$1/0.02 = 50$. With $F$ between 10 and 100, report the ITT and first stage, plus an Anderson--Rubin set; do not headline the ratio.}

\begin{exercise}[judgement]
An analyst proposes using rainfall as an instrument for QuickBite's delivery fee (rain reduces courier supply, which raises the
surge fee) to estimate the fee elasticity of orders. Assess each of the four assumptions.
\end{exercise}
\answer{Relevance: checkable, likely strong. Independence: rain is not chosen by anyone, but it is seasonal and regional; compare within
city and season. Monotonicity: rain raises fees everywhere, plausibly. Exclusion: almost certainly fails, because rain raises demand for
delivery directly (people stay in), so the ratio mixes the fee effect with the rain effect on demand, biasing the elasticity upward
(towards zero or positive). A better instrument shifts the fee without shifting demand, such as randomised fee tests or courier-pay
changes.}

%=====================================================================================
\sectionthree

\begin{extension}[Continuous treatments, flexible controls and describing compliers]
With a continuous treatment and binary instrument, the Wald ratio is a weighted average of unit slopes over the range the instrument
moved them (Angrist, Graddy and Imbens, 2000). With many controls, residualise $Y$, $D$ and $Z$ on $X$ with cross-fitted learners and compute
$\sum\tilde Z\tilde Y/\sum\tilde Z\tilde D$: the partially linear IV model of double machine learning. Abadie's $\kappa$-weighting identifies any feature of the
compliers' joint distribution.
\end{extension}

\begin{extension}[Sensitivity to exclusion]
Report the LATE as a function of the direct effect, $(\text{ITT} - \delta)/\pi$, over a range of $\delta$ argued from the institution, and the
$\delta$ at which the decision flips. Conley, Hansen and Rossi (2012) formalise this as ``plausibly exogenous'' instruments.
\end{extension}

\subsection*{Annotated reading}
\begin{readinglist}
\reading{Angrist, Imbens and Rubin (1996), ``Identification of causal effects using instrumental variables'', \emph{JASA}}{The
potential-outcomes framework for IV, compliance types and the LATE}{Sections 1--4}{2}
\reading{Angrist and Pischke (2015), \emph{Mastering 'Metrics}, Princeton}{IV at an intuitive level, with classic
examples}{Ch.\ 3}{1}
\reading{Imbens and Angrist (1994), ``Identification and estimation of local average treatment effects'',
\emph{Econometrica}}{The LATE theorem}{All}{3}
\reading{Lee, McCrary, Moreira and Porter (2022), ``Valid t-ratio inference for IV'', \emph{AER}}{Why $F > 10$ is not enough, and
corrected critical values}{Sections I--III}{3}
\reading{Andrews, Stock and Sun (2019), ``Weak instruments in instrumental variables regression'', \emph{Annual Review of
Economics}}{A survey of weak-instrument-robust inference}{Sections 1--4}{3}
\reading{Graddy (2006), ``The Fulton Fish Market'', \emph{Journal of Economic Perspectives}}{Weather as an instrument for price in a
real market}{All}{1}
\end{readinglist}

\end{document}
